% !PRINTED 74
\weberpage{74}

%<<CONT
in which

\[
\varepsilon_0 = e^{\frac{\pi i \omega}{4}}
\]

and

\weq{11}{
\vartheta_{11}(0) = 0,\quad \vartheta_{10}\!\left(\tfrac{1}{2}\right) = 0,\quad \vartheta_{01}\!\left(\tfrac{\omega}{2}\right) = 0,\quad \vartheta_{00}\!\left(\tfrac{1+\omega}{2}\right) = 0.
}

\gesperrt{The squares of the functions (7):}

\weq{12}{
\vartheta_{00}^2(v),\quad \vartheta_{01}^2(v),\quad \vartheta_{10}^2(v),\quad \vartheta_{11}^2(v)
}

are $\Theta_{00}$-functions of the second order, and consequently there
exist two linear relations between them. If we take these relations
 in the form:

\[
\begin{aligned}
\vartheta_{10}^2(v) &= A\,\vartheta_{01}^2(v) + B\,\vartheta_{11}^2(v),\\
\vartheta_{00}^2(v) &= A'\vartheta_{01}^2(v) + B'\vartheta_{11}^2(v),
\end{aligned}
\]

then one can easily determine the coefficients on the basis of formulae (10), (11),
by setting $v = 0$ and $v = \dfrac{\omega}{2}$. In this way one
obtains

\weq{13}{
\begin{aligned}
\vartheta_{01}^2\,\vartheta_{10}^2(v) &= \vartheta_{10}^2\,\vartheta_{01}^2(v) - \vartheta_{00}^2\,\vartheta_{11}^2(v),\\
\vartheta_{01}^2\,\vartheta_{00}^2(v) &= \vartheta_{00}^2\,\vartheta_{01}^2(v) - \vartheta_{10}^2\,\vartheta_{11}^2(v),
\end{aligned}
}

and from this further, on setting $v = \tfrac{1}{2}$:

\weq{14}{
\vartheta_{00}^4 = \vartheta_{01}^4 + \vartheta_{10}^4.
}

By means of any two of the squares (12) all functions $\Theta_{00}$ of
the second order can be represented linearly; but the remaining
$\Theta$-functions of the second order can also be composed from the
functions $\vartheta_{g_1,g_2}$, for one obtains for each
characteristic two linearly independent products, of which the one is
an even and the other an odd function, namely:

\weq{15}{
\begin{aligned}
\vartheta_{00}(v)\,\vartheta_{01}(v), &\quad \vartheta_{10}(v)\,\vartheta_{11}(v) & &\text{Characteristic } (0,1)\\
\vartheta_{00}(v)\,\vartheta_{10}(v), &\quad \vartheta_{01}(v)\,\vartheta_{11}(v) & &\text{„}\quad (1,0)\\
\vartheta_{10}(v)\,\vartheta_{01}(v), &\quad \vartheta_{00}(v)\,\vartheta_{11}(v) & &\text{„}\quad (1,1)
\end{aligned}
}

According to the same principle all $\Theta$-functions of arbitrary
order and of arbitrary principal characteristic can now be formed from
the functions $\vartheta$. To establish this, we denote by
$\Theta_0, \Theta_1$ any two of the $\vartheta$-squares (12), and by $F^{(n)}(\Theta_0,\Theta_1)$ an
integral rational and homogeneous function of order $n$ of the two
arguments $\Theta_0, \Theta_1$. In accordance with this one obtains for the $\Theta$-functions
of order $m$, $\Theta^{(m)}(v)$, the following expressions:

%CONT>>
